% !TEX program = xelatex

\documentclass{cv}

% 关闭页码显示
\pagenumbering{gobble}

% 启用中文支持
\useChinese{Regular}{Bold}

\begin{document}

%% 基本信息（公开版：不放照片、出生年月、考研成绩等个人信息）
\name{王子隆}
\begin{basicinfo}
  \info{\email{zilong\_@stu.xjtu.edu.cn}}
  \info{\phone{. }\wechat{15534516801}}
  \info{\github{xiuxiua8}[https://github.com/xiuxiua8]}
  \info{\linkedinsquare{zilong-wang-xiu}[https://www.linkedin.com/in/zilong-wang-xiu/]}
\end{basicinfo}

\section{教育经历}
\cventry{西安交通大学}
  {2026年9月 -- 至今}
  [计算机技术 | 专业硕士]
  [西安]

\cventry{西安交通大学}
  {2022年9月 -- 2026年6月}
  [计算机科学与技术 | 工学学士]
  [西安]
  \begin{itemize}
    \item \textbf{成绩}：平均分 84.5/100
    \item \textbf{核心课程}：操作系统(A)、算法分析与设计(A)、数据结构(A)、计算机系统导论(A+)、凸优化(A+)
  \end{itemize}

\section{专业实习}
\cventry{西安交通大学 | 校内暑期实习}
  {2025年7月}
  [PromptMeet 智能会议纪要系统 | 项目组长]
  [\iconlink[\faGithub][PromptMeet]{https://github.com/xiuxiua8/PromptMeet}]
  \begin{itemize}
    \item 担任10人团队组长，负责需求澄清与架构规划，将系统拆解为转录、纪要生成、任务提取三大模块
    \item 主导选型 FastAPI + Vue.js + MySQL，通过 WebSocket 与异步服务实现前后端解耦，支撑多人并行开发
    \item 基于 Whisper 完成中文实时转录，结合 DeepSeek 大模型自动输出发言摘要、行动项等结构化纪要
  \end{itemize}

\section{项目经历}
\cventry{PromptMeet macOS 原生端}
  {2026年7月 -- 2026年8月}
  [个人迭代 | Swift / SwiftUI / ScreenCaptureKit / whisper.cpp]
  [\iconlink[\faGithub][PromptMeet]{https://github.com/xiuxiua8/PromptMeet}]
  \begin{itemize}
    \item 以 ScreenCaptureKit 与 AVAudioEngine 双路采集会议声音和麦克风，转写时区分“我”与“会议”
    \item 默认以本机 whisper.cpp 实时转写，会议音频不上传云端；任一音源不可用时另一路继续工作
    \item 按会议持久化转写、截图分析、AI 问答、摘要与待办的时间线，支持历史会议回看与继续提问
  \end{itemize}

\cventry{AgentKeyLight 键盘状态灯}
  {2026年9月}
  [个人项目 | Python / macOS / USB HID]
  [\iconlink[\faGithub][AgentKeyLight]{https://github.com/xiuxiua8/AgentKeyLight}]
  \begin{itemize}
    \item 通过 hooks 汇总 Codex 与 Claude Code 多会话状态，以逐键背光动画显示运行、等待、完成、出错
    \item 通过 USB HID 推送临时颜色帧，不改写键盘内保存的灯效；处理多会话优先级、进程退出与状态过期
    \item 提供本机网页调色板与诊断命令，支持实时预览、颜色校正与配置热加载
  \end{itemize}

\cventry{计算机系统导论 (CSAPP)}
  {\iconlink[\faGithub][ics]{https://github.com/xiuxiua8/ics}}
  [计算机系统 | C / ELF / GDB]
  []
  \begin{itemize}
    \item 自学 CMU 15-213，覆盖体系结构、缓存优化、内存管理、链接与程序安全等核心主题
    \item 实现静态链接器，梳理 ELF 重定位流程，加深对程序加载链路的理解
    \item 进行缓冲区溢出与 ROP 实验，分析栈帧并利用 GDB 构造输入，验证控制流劫持与防护手段
  \end{itemize}

\cventry{区块链系统}
  {\iconlink[\faGithub][blockchainRUST]{https://github.com/xiuxiua8/blockchainRUST}}
  [分布式系统 | Rust / Tokio / libp2p]
  []
  \begin{itemize}
    \item 使用 Rust 实现区块链系统原型，包括区块结构、工作量证明、UTXO 交易、钱包签名与链状态持久化
    \item 基于 libp2p 搭建节点间网络通信，支持交易广播与区块同步，并编写集成测试验证核心功能
  \end{itemize}

\cventry{CS61B / CS188 课程项目}
  {\iconlink[\faGithub][cs61b]{https://github.com/xiuxiua8/CS61B}}
  [数据结构与算法 | 人工智能]
  []
  \begin{itemize}
    \item 自学 UC Berkeley CS61B，实现图搜索、排序、并查集等数据结构与算法，代码通过课程自动评测
    \item 自学 UC Berkeley CS188，在 Pac-Man 项目中实现 A* 搜索、Alpha-Beta 剪枝与强化学习算法
  \end{itemize}

\section{竞赛经历}
\cventry{大学生数学建模竞赛}
  {\iconlink[\faGithub][solution]{https://github.com/xiuxiua8/24_problem_b}}
  [数学建模]
  [校级二等奖]
  \begin{itemize}
    \item 负责代码与建模，基于 K-Means 聚类和卡方检验研究心理测试数据分类
  \end{itemize}

\section{技能证书}
\begin{compactlist}
  \item \textbf{编程语言}：熟练 C++ / Rust / Java / Python，使用 Swift 开发 macOS 应用，习惯用自动化测试保障改动
  \item \textbf{工程实践}：FastAPI、Vue.js、SwiftUI、WebSocket、MySQL、Git；了解 P2P 网络、共识算法与异步架构
  \item \textbf{语言成绩}：CET-4 548分，CET-6 550分
\end{compactlist}

\end{document}
